\lyseis{Β3.13}{Ευθεία και κύκλος}
\sol{1} Για $d=4<6$: τέμνουσα, $2$ κοινά σημεία. Για $d=6=6$: εφαπτομένη, $1$ κοινό σημείο. Για $d=9>6$: εξωτερική, κανένα κοινό σημείο.
\sol{2} α) $3<8$: δύο. β) $4=4$: ένα. γ) $6>3$: κανένα. δ) $2{,}5=2{,}5$: ένα. ε) $12>10$: κανένα.
\sol{3} α) Σωστό. β) Σωστό. γ) Λάθος, είναι μεγαλύτερη. δ) Σωστό. ε) Λάθος, το πολύ δύο. στ) Σωστό (απόσταση $0<\rho$).
\sol{4} Ελεύθερη κατασκευή. Η ευθεία με $3<4$ είναι τέμνουσα, με $4=4$ εφαπτομένη, με $5>4$ εξωτερική.
\sol{5} Ελεύθερη κατασκευή. Φέρνεις την ακτίνα $\mathrm{ΟΑ}$ και μετά την κάθετη στην $\mathrm{ΟΑ}$ στο $\mathrm{Α}$.
\sol{6} Ένα κοινό σημείο (η $\varepsilon$ εφάπτεται). Για δύο κοινά σημεία πρέπει η απόσταση να γίνει μικρότερη από $3$ cm.
\sol{7} Στα $2{,}5$ m ($2{,}5<4$) το μονοπάτι τέμνει τον κύκλο και βρέχονται δύο σημεία (και ό,τι ανάμεσα τους). Στα $4$ m ($4=4$) το μονοπάτι εφάπτεται και βρέχεται ένα μόνο σημείο. Στα $6$ m ($6>4$) δεν βρέχεται κανένα.
\sol{8} Ο τροχός εφάπτεται στον δρόμο, άρα ακουμπά σε ένα σημείο. Ο άξονας απέχει από τον δρόμο όσο η ακτίνα, $35$ cm.
\sol{9} Το κέντρο απέχει $12$ cm και από τον τοίχο και από το δάπεδο, γιατί η μπάλα εφάπτεται και στα δύο (απόσταση ίση με την ακτίνα). Ο τοίχος και το δάπεδο είναι δύο εφαπτόμενες ευθείες του κύκλου.
\sol{10} Όχι. Η απόσταση $38<50$, άρα ο δρόμος θα ήταν τέμνουσα και θα διέσχιζε τη λίμνη. Για να περνά στην άκρη χωρίς να τη διασχίζει, η απόσταση πρέπει να είναι $50$ m (εφαπτομένη) ή περισσότερο.
\sol{11} Στα $30$ m ($30>25$) το σήμα δεν φτάνει στον δρόμο. Στα $25$ m ($25=25$) φτάνει σε ένα μόνο σημείο. Στα $20$ m ($20<25$) το σήμα φτάνει σε δύο σημεία και σε ολόκληρο το τμήμα του δρόμου ανάμεσά τους.
\sol{12} Οι εφαπτόμενες στα $\mathrm{Α}$ και $\mathrm{Β}$ είναι κάθετες στην $\mathrm{ΑΒ}$, άρα παράλληλες μεταξύ τους, και απέχουν $\mathrm{ΑΒ}=2\rho$. Το ίδιο ισχύει για τις εφαπτόμενες στα $\mathrm{Γ}$, $\mathrm{Δ}$, που είναι κάθετες στην $\mathrm{ΓΔ}$. Επειδή $\mathrm{ΑΒ}\perp\mathrm{ΓΔ}$, οι δύο ομάδες είναι κάθετες μεταξύ τους. Σχηματίζεται τετράγωνο με πλευρά $2\rho$.
\sol{13} Ελεύθερη κατασκευή. Φέρνεις τη διάμετρο $\mathrm{ΓΔ}$ κάθετη στην $\mathrm{ΑΒ}$ και τις εφαπτόμενες στα $\mathrm{Γ}$, $\mathrm{Δ}$. Και οι δύο είναι κάθετες στην $\mathrm{ΓΔ}$, όπως και η $\mathrm{ΑΒ}$, άρα παράλληλες προς την $\mathrm{ΑΒ}$. Είναι δύο.
\sol{14} Στα σημεία επαφής $\mathrm{Α}$ και $\mathrm{Β}$ φέρνουμε τις κάθετες στις $\varepsilon_1$ και $\varepsilon_2$ αντίστοιχα. Η ακτίνα είναι κάθετη στην εφαπτομένη, άρα και οι δύο κάθετες περνούν από το κέντρο. Το κέντρο είναι το σημείο τομής τους.
\sol{15} α) Μία. β) Δύο (φαίνεται με χάρακα). γ) Καμία: κάθε ευθεία που περνά από εσωτερικό σημείο έχει απόσταση από το κέντρο μικρότερη από την ακτίνα, άρα είναι τέμνουσα.
\sol{16} Ελεύθερη εργασία. Η απόσταση του κέντρου από τον χάρακα πρέπει να βγαίνει ίση με την ακτίνα του καπακιού.
\sol{17} Ελεύθερη κατασκευή. Και τα δύο κέντρα απέχουν από τον δρόμο όσο η ακτίνα, δηλαδή βρίσκονται στην ίδια απόσταση από αυτόν, πάνω σε ευθεία παράλληλη προς τον δρόμο.
